
\documentclass[11pt,a4paper]{article}
\usepackage[utf8]{inputenc}
\usepackage[T1]{fontenc}
\usepackage{lmodern}
\usepackage{amsmath,amssymb,amsthm,mathtools}
\usepackage{geometry,enumitem,xcolor,hyperref,booktabs}
\geometry{top=2.5cm,bottom=2.5cm,left=2.5cm,right=2.5cm}
\hypersetup{colorlinks=true,linkcolor=blue!60!black,urlcolor=blue!60!black,citecolor=blue!60!black}
\newtheorem{theorem}{Theorem}[section]
\newtheorem{lemma}[theorem]{Lemma}
\newtheorem{proposition}[theorem]{Proposition}
\newtheorem{definition}[theorem]{Definition}
\newtheorem{axiom}[theorem]{Axiom}
\newtheorem{corollary}[theorem]{Corollary}
\newtheorem{remark}[theorem]{Remark}
\newcommand{\WCC}{\mathcal{W}}
\newcommand{\PWCC}{\mathbf{P}_{\mathrm{WCC}}}
\newcommand{\NPWCC}{\mathbf{NP}_{\mathrm{WCC}}}
\newcommand{\PTuring}{\mathbf{P}_{\mathrm{Turing}}}
\newcommand{\NPTuring}{\mathbf{NP}_{\mathrm{Turing}}}
\newcommand{\CTheta}{C_{\Theta,\varepsilon}}
\newcommand{\poly}{\operatorname{poly}}
\title{\textbf{Classical Complexity and Curvature-Bounded Physical Computation}\\[4pt]
\large A Critical Reassessment of the Proposed Physical $P$ versus $NP$ Separation}
\author{Richard J. Reyes}
\date{Original research note: December 1, 2025\\Revised mathematical audit: October 27, 2026}
\begin{document}
\maketitle
\begin{abstract}
This technical revision distinguishes the mathematically well-defined classical
$P$ versus $NP$ question from questions of realizability and resource usage
in Wave Curvature Computation (WCC). The discrete finite-state, local-update
model can represent polynomial-time Turing computations under uniform encoding,
decoding, and transition assumptions. That equivalence does not prove that each
symbolic computation has a polynomial-curvature physical realization.
Several claims in the original December 2025 technical note required correction:
physical distinguishability alone does not force positive Laplacian-curvature
cost; finite first-derivative energy and locality do not imply a universal
polynomial bound on second-derivative curvature; expander graphs do not, without
a separate geometric embedding theorem, imply exponential interface area; and
the total influence of any $n$-bit Boolean function is at most $n$.
We establish exact counterexamples and corrected conditional estimates,
clarify physical $P$ and $NP$ witness admissibility, and isolate explicit proof
obligations for a possible restricted-model separation. No unconditional
separation of classical $P$ and $NP$, nor a proved separation of the
curvature-bounded physical classes, is claimed here.
\end{abstract}
\paragraph{New decoder-uniformity obstruction (November 2, 2026).}
A fixed, language-independent sign map encodes every $n$-bit string into
$n$ real amplitudes $\pm1$, with squared mass $n$, and the sign receiver
recovers the string under per-coordinate additive errors with magnitude
strictly below one. Composing this receiver with an arbitrary predicate
would decide that predicate without charging the final decoding step.
Such a decoder need not be computable, let alone efficient.
Consequently, $2^n$ recoverable labels, finite amplitude cost,
and a uniform encoding cannot establish an unconditional lower bound
for decision problems without bounding the decoder's complexity.
This counterexample does not establish an efficient SAT algorithm.
The physical continuum-to-computation bridge and universal
curvature-versus-decision lower bound remain open.



\section{Historical Context: The Divergence of Logic and Physics}
The classical $P$ versus $NP$ question concerns language classes defined by
uniform Turing-machine running time and polynomial verification; the question
is mathematically well posed independent of a physical implementation
\cite{Sipser,AroraBarak2009}. Turing machines are abstract algorithms and do
not assert that unbounded physical memory is available at zero energetic cost.
Whether a concrete physical computer efficiently realizes a particular
algorithm is an additional question. Locality, thermodynamic dissipation,
finite precision, field regularity, and resource budgets can restrict a
physical substrate without changing the definition of classical $P$ or $NP$.

WCT motivates such a resource-sensitive investigation. We distinguish three
notions throughout: classical classes $P,NP$; uniformly specified discrete
WCC classes $P^{\mathrm{disc}}_{\mathrm{WCC}},NP^{\mathrm{disc}}_{\mathrm{WCC}}$;
and physical, curvature-rail classes
$P^{\mathrm{phys}}_{\mathrm{WCC}},NP^{\mathrm{phys}}_{\mathrm{WCC}}$.
A physical-vs-classical comparison must specify a uniform implementation
map and its costs before inferring equality or separation of these classes.

\section{The Physical Distinguishability Obstruction}
\label{sec:distinguishability}
\subsection{Physical Encoding of Logical Bits}
Let $\Omega\subset\mathbb R^3$ be a bounded Lipschitz domain, with
periodic or Neumann boundary conditions allowed. Two field encodings
$\psi_0,\psi_1\in H^2(\Omega)$ can be distinct in the $H^1$ norm without
having different Laplacians. We therefore distinguish \emph{symbolic
separation} from \emph{curvature cost}.

\begin{definition}[Regularized curvature cost]
\label{def:regularized}
For $\varepsilon>0$ and $\psi\in H^2(\Omega;\mathbb C)$ define the
complex-safe curvature operator
\[
\Theta_\varepsilon[\psi]
 =-\frac{\Delta\psi\,\overline\psi}{|\psi|^2+\varepsilon^2},
\qquad
\CTheta[\psi]=\int_\Omega|\Theta_\varepsilon[\psi]|^2\,dx.
\]
For a pair of encodings, a proposed separation cost is
$\CTheta[\psi_1-\psi_0]$. This particular functional measures
second-derivative structure, not every possible resource cost of a bit.
\end{definition}

\begin{proposition}[Counterexample to a universal positive separation cost]
\label{prop:zero_cost}
Distinct $H^1$ states need not have positive curvature-separation cost.
\end{proposition}
\begin{proof}
Set $\psi_0\equiv0$ and $\psi_1\equiv1$ on $\Omega$. Then
$\|\psi_1-\psi_0\|_{H^1}^2=|\Omega|>0$ while
$\Delta(\psi_1-\psi_0)=0$. Hence
\[
\CTheta[\psi_1-\psi_0]=0.
\]
The counterexample obeys the stated periodic or Neumann conditions.
\end{proof}

\begin{remark}[WIN-12]
The original Lemma 2.2 inferred nonzero Laplacian from nonzero $H^1$
distance, which is false. An experimentally grounded lower bound for a
specific device would require extra boundary conditions, state constraints,
and a coercivity inequality; distinguishability alone is insufficient.
\end{remark}

\subsection{Curvature Cost of an \texorpdfstring{$n$}{n}-Bit Configuration}
A lower bound $c_0n$ for storing $n$ bits may be established within a
specified \emph{independent-cell encoding} if every occupied cell contributes
at least $c_0>0$ and the chosen functional is additive across the disjoint
supports. Neither cell independence nor positive $c_0$ follows from the
previous definition. In particular, an $n$-bit string is not required to
simultaneously realize the $2^n$ possible assignments.

\begin{proposition}[Conditional additive storage estimate]
If $\Omega=\bigsqcup_{j=1}^n\Omega_j$ and a chosen nonnegative storage
functional $E_{\rm store}$ is additive on the disjoint cells with
$E_{\rm store}(\psi_{b_j}|_{\Omega_j})\ge c_0$ for each $j$, then
$E_{\rm store}(\psi_w)\ge nc_0$.
\end{proposition}
\begin{proof}
Sum the $n$ cellwise inequalities. No such conclusion is asserted for
$\CTheta$ without proving the premises.
\end{proof}

\subsection{Classical Turing Machines and Physical Implementations}
The Turing model accounts for operations through a formal step count, not
through a PDE-based curvature-energy functional. Its omission of an
explicit energy budget does not imply a theorem that every physical
implementation must consume zero energy or that a polynomial-overhead
physical simulation is impossible. The physical cost of symbols, memory,
and erasure is an independent quantitative engineering question
\cite{Landauer1961,Bennett1982}.

\subsection{Physical Complexity Classes}
For a fixed, uniform physical model $\mathfrak M$ and a stated polynomial
resource rail $R$, define $P^{\mathrm{phys}}_{\mathfrak M,R}$ by uniformly
specified deterministic deciders that meet the rail on every input. Define
$NP^{\mathrm{phys}}_{\mathfrak M,R}$ with polynomial-length, physically
admissible witness encodings and deterministic verifiers subject to the
same rail. The physical classes are machine-relative; we do not identify
them with classical classes without separately proving the necessary
simulation theorems.

\section{Classical Computation and Physical Resource Models}
\label{sec:tm_not_physical}
The classical formalism is mathematically defined even if no hardware can
realize an unbounded tape for all input sizes at once. Complexity theory
uses a \emph{family} of finite computations, one for each input length.
Physical realizability concerns whether a single uniform device family has
allowed costs, scaling, and error control. The extended Church--Turing
principle is a statement about efficient simulation of physical models, not
a prerequisite for formulating classical $P$ versus $NP$.

\begin{remark}
A distinction between two physical encoding conventions cannot by itself
refute a theorem about abstract languages. Conversely, a formal equivalence
of symbolic update systems does not guarantee faithful realization by a
specified nonlinear PDE under fixed energy and stability rails.
\end{remark}

\section{The WCC Model and Curvature Axioms}
\label{sec:wcc_axioms}
\begin{definition}[Polynomial curvature rail]
For a uniform physical device family indexed by input length $n$, a
trajectory $\psi_{n,x}(t)$ respects a polynomial curvature rail if there
exists one polynomial $p$ such that, for every input $x$ of length $n$ and
all times $0\le t\le T(n)$,
\[
\sup_{|x|=n}\sup_{0\le t\le T(n)}
\CTheta[\psi_{n,x}(t)]\le p(n).
\]
A requirement of \emph{finite} curvature at each $n$ is strictly weaker.
\end{definition}

\begin{lemma}[A valid conditional curvature bound]
\label{lem:h2bound}
For every $\psi\in H^2(\Omega;\mathbb C)$ and $\varepsilon>0$,
\[
\CTheta[\psi]\le\frac{1}{4\varepsilon^2}
\|\Delta\psi\|_{L^2(\Omega)}^2.
\]
\end{lemma}
\begin{proof}
For $r\ge0$, $r/(r^2+\varepsilon^2)\le1/(2\varepsilon)$.
Apply this pointwise to $r=|\psi|$ and integrate the square.
\end{proof}

\begin{remark}[WIN-13]
The original operator $-\Delta\psi/(\psi+\varepsilon)$ does not have the
claimed uniform denominator bound $|\psi+\varepsilon|\ge\varepsilon$:
$\psi=-\varepsilon$ is an immediate counterexample. One may instead
assume an explicit gap $|\psi+\varepsilon|\ge\delta>0$ or use the
complex-safe regularizer of Definition~\ref{def:regularized}. All estimates
in this revision use the latter.
\end{remark}

\section{Discrete and Physical WCC Semantics}
\label{sec:two_semantics}
\subsection{Discrete WCC Model}
A discrete WCC machine uses a finite alphabet $S$, a fixed-radius local
update rule, polynomially many lattice cells, a uniform family of lattices,
and polynomial-time constructible encoders and decoders. Machine
descriptions must not contain nonuniform advice or a hard language in the
initialization map.

\begin{theorem}[Discrete Turing equivalence under uniformity]
\label{thm:disc-equiv}
For the uniform finite-state, fixed-radius WCC model defined above, with
polynomial configuration size and polynomial rounds,
\[
P^{\mathrm{disc}}_{\mathrm{WCC}}=P,
\qquad NP^{\mathrm{disc}}_{\mathrm{WCC}}=NP.
\]
\end{theorem}
\begin{proof}
A Turing machine sequentially updates all polynomially many cells in
polynomial time per global WCC step. Conversely, standard finite-state
local tape-and-head encoding simulates each step of a fixed polynomial-time
Turing machine by a constant number of local rounds on a polynomial-size
lattice. The two simulations preserve polynomial witness length and verifier
runtime. Polynomial-time uniform encoding, decoding, and lattice generation
exclude hidden oracle computations.
\end{proof}

\subsection{Physical WCC Model and Witness Admissibility}
The physical model additionally restricts allowable embeddings, precision,
curvature-energy budgets, and dynamics. Therefore the discrete theorem
alone supplies no equality
$P^{\mathrm{phys}}_{\mathrm{WCC}}=P$ or
$NP^{\mathrm{phys}}_{\mathrm{WCC}}=NP$.

\begin{theorem}[Verifier inclusion; correction of the empty-NP claim]
\label{thm:p-in-np}
In any model whose NP verifiers may ignore their witnesses and which
subjects deciders and verifiers to the same resource rails,
\[
P^{\mathrm{phys}}_{\mathrm{WCC}}
\subseteq NP^{\mathrm{phys}}_{\mathrm{WCC}}.
\]
\end{theorem}
\begin{proof}
Given a polynomial-resource decider $D(x)$, define
$V(x,w)=D(x)$ for a fixed empty witness (or ignore $w$). It has the
same resource usage and correctness. Thus each physical-P language has
a physical-NP verifier.
\end{proof}
\begin{remark}[WIN-16]
The original assertion $NP^{\mathrm{phys}}_{\mathrm{WCC}}=\varnothing$
contradicts the claimed nonempty physical-P class together with
Theorem~\ref{thm:p-in-np}. No physical class separation follows from
that assertion.
\end{remark}

\subsection{A Correct Comparison Statement}
\begin{proposition}[No logical conflict between the two semantics]
Suppose that the discrete classes equal classical $P,NP$, but physical
admissibility imposes additional restrictions. It is logically possible
that the physical classes are proper subsets of their respective discrete
classes. No equality or inequality between physical P and NP follows from
those facts alone.
\end{proposition}
\begin{proof}
A class inclusion into a larger class does not determine whether two
subclasses are equal. In particular the implication
$P^{\mathrm{phys}}\subseteq P^{\mathrm{disc}}$ and
$NP^{\mathrm{phys}}\subseteq NP^{\mathrm{disc}}$ is compatible with both
$P^{\mathrm{phys}}=NP^{\mathrm{phys}}$ and a strict inclusion.
\end{proof}

\subsection{On the Claimed Ill-Posedness}
\begin{remark}[WIN-17]
Classical $P$ versus $NP$ is a precise question about languages and
Turing machines. What can be \emph{underspecified} is the expression
``physically efficient computation'' unless the physical machine,
precision, allowed resources, and asymptotic family are specified.
Failure of a physical equivalence theorem cannot make the classical
mathematical question ill-posed.
\end{remark}

\section{Locked Expanders and Physical Embedding Costs}
\label{sec:locked}
A bounded-degree expander has an edge boundary of order the size of a
balanced vertex subset. That is an $\Omega(n)$ combinatorial cut
constraint on $O(n)$ vertices; it does not by itself prove that \emph{one}
encoded satisfying assignment must have $2^{\Omega(n)}$ interface area.
Counting exponentially many assignments is different from physically
storing all assignments at once.

\begin{proposition}[Conditional interface-to-curvature lower bound]
\label{prop:interface}
Suppose that for each admissible realization $\psi$ of an instance $F_n$
there is a geometric interface $\Gamma(\psi)$ such that
\[
\mathcal H^2(\Gamma(\psi))\ge A_n
\quad\text{and}\quad
\CTheta[\psi]\ge\kappa\,\mathcal H^2(\Gamma(\psi))
\]
with a constant $\kappa>0$ uniform in $n$, the input, and the encoding.
Then every such realization satisfies $\CTheta[\psi]\ge\kappa A_n$.
\end{proposition}
\begin{proof}
Compose the two assumed inequalities.
\end{proof}
\begin{remark}[WIN-18]
A proved $A_n\ge 2^{\beta n}$, uniformly for every admissible
implementation of a specific \emph{decision problem}, would be a
substantive new embedding theorem. Graph expansion alone does not
establish it, and the original exponential-domain-wall theorem is
therefore withdrawn as an unconditional conclusion.
\end{remark}

\begin{proposition}[Verification of a supplied assignment]
\label{prop:assignment-P}
The language
\[
L_{\mathrm{check}}=\{(F,\alpha):\alpha\text{ satisfies the 3-CNF formula }F\}
\]
is in classical $P$.
\end{proposition}
\begin{proof}
Evaluate each literal of each clause under the supplied assignment and
accept exactly when every clause contains a true literal. The runtime is
polynomial in the encoded formula and assignment lengths. Restricting to
uniformly constructible expander formulas does not make this
verification task NP-hard.
\end{proof}

\section{A Conditional Constructibility Barrier}
\label{sec:constructibility}
\begin{definition}[Per-step curvature increment]
For a trajectory $\psi_{t+1}=F_n(\psi_t)$ set
$\Delta\CTheta(t)=\CTheta[\psi_{t+1}]-\CTheta[\psi_t]$.
A bound $|\Delta\CTheta(t)|\le K(n)$ is a \emph{separate hypothesis};
finite precision and local update radius alone do not imply a uniform
constant $K$ for all system sizes.
\end{definition}

\begin{theorem}[Conditional telescoping capacity estimate]
\label{thm:constructibility}
If $\CTheta[\psi_0]\le p_0(n)$,
$|\Delta\CTheta(t)|\le K(n)$ for all $t<T(n)$,
$T(n)\le p_T(n)$, and $p_0,K,p_T$ are polynomials, then
\[
\CTheta[\psi_{T(n)}]\le p_0(n)+p_T(n)K(n)=\poly(n).
\]
\end{theorem}
\begin{proof}
Telescope the increments and use the triangle inequality:
\[
\CTheta[\psi_T]=\CTheta[\psi_0]+\sum_{t<T}\Delta\CTheta(t)
\le p_0(n)+T(n)K(n).
\]
\end{proof}
\begin{remark}[WIN-19]
A high-curvature \emph{witness} that violates the physical verifier rail
is not evidence that its language belongs to physical NP. A separation
requires one common language $L$ with a polynomial-resource physical
verifier for every admissible witness and a proof that no permitted
polynomial-resource physical decider recognizes $L$. A lower bound only
on one chosen witness representation is insufficient.
\end{remark}

\section{Turing Simulability of Discrete WCC}
The discrete simulation of Theorem~\ref{thm:disc-equiv} is an algebraic
statement about locally updated finite symbols. With $N(n)$ cells and
$T(n)$ global rounds, a sequential machine can evaluate the entire
evolution in $T(n)\poly(N(n))$ time. This does not convert an arbitrary
physical PDE into the desired finite-state machine, nor does it charge
curvature resources to the simulation. Such a cost-preserving realization
must be established independently (see \cite{Reyes2025_DiscreteWCC}).

\section{Classical-to-Physical Implications}
\label{sec:ECT}
The extended Church--Turing thesis addresses simulation of physical
computations in an abstract model with polynomial overhead. A
one-directional simulation such as
$P^{\mathrm{phys}}\subseteq P$ does not entail
$P^{\mathrm{phys}}=P$, nor does it identify physical NP with classical NP.
A physical equality claim would require the relevant simulation in both
directions and preservation of witnesses, encodings, errors, and budgets.

\begin{remark}
Even the hypothetical equality $P=NP$ would not, solely by definition,
establish an equality between two separately defined physical classes.
Conversely, unequal physical subclasses would not by themselves resolve
classical $P$ versus $NP$. No ``if and only if ECT'' criterion is proved
or assumed here.
\end{remark}

\paragraph{Updated scope of current computational results.}
The repository contains a uniform Cook--Levin
reduction and restricted quadratic-type single-tape
SAT lower bounds.  Neither a superpolynomial lower
bound for every correct SAT decider nor a
uniform physical-to-Turing simulation under the
specified WCT resource constraints has been proved.
Formula-dependent phase-energy certificates do
not provide efficient physical preparation of the
required field.  The newer results therefore do
not overturn the conditional physical conclusions
of this reassessment.

\section{Conclusion: What Is Established and What Remains Open}
\label{sec:conclusion}
The corrected statements are: (i) under explicit uniformity, discrete WCC
has the standard polynomial Turing simulation and verifier equivalence;
(ii) positive storage or curvature costs require device-specific coercivity
assumptions; (iii) $H^2$ control together with a positive regularization
parameter gives a rigorous curvature upper bound; (iv) expander
connectivity does not independently establish an exponential lower bound
on the curvature of every decision procedure; and (v) Boolean influence
cannot provide an exponential lower bound because it never exceeds $n$.

The proposed physical separation remains an open mathematical program.
A viable theorem must name a language, specify its uniform physical
encoding and resource rails, provide an admissible polynomial-resource
physical verifier, and prove a lower bound against every admissible
physical decider, not only against a particular geometric search protocol.
Classical $P$ versus $NP$ remains mathematically well posed and unresolved
by these constructions.

\subsection{Current Lean formalization status (7 November 2026)}
\label{sec:current-lean-status}
The revised Part III inventory contains 22 entries:
21 Exact under their explicit hypotheses and one
Open interpretive remark (III-22), concerning the
arrow of time.  Exact includes \emph{conditional}
analytic results whose physical assumptions
are not independent Lean theorems.
This unified register supersedes
the separate dated verification addenda.
\emph{Addendum, 13 November 2026.}  III-22 is now classified as
Interpretive rather than Open: the arrow-of-time remark states that no
equivalence is proved and asserts no theorem, so no proof obligation is
attached to it.

\paragraph{Analytical and finite mathematical results.}
The denominator counterexample, the zero
symbolic-curvature example, the
H$^2$-to-curvature comparison
(Lemma~\ref{lem:h2bound}), conditional
storage and channel-capacity estimates
(Theorem~\ref{thm:capacity}), and the
high-frequency periodic-cube example
(Proposition~\ref{prop:highfreq}) have
matching scoped Lean results.
The Boolean-hypercube influence bounds
and the supplied CNF-assignment
verifier are also formalized.  The
supplied assignment test lies in
polynomial time and is not a
candidate separation language.

\paragraph{Discrete computational equivalence.}
The finite-state, fixed-radius discrete
models have explicit Turing simulations,
including all-dimensional sweep models
with charged costs.
The actual uniform lattice model,
including polynomial-time lattice
generation, encoding and readout,
has exact theorems
\texttt{PartIIUniform.P\_disc\_eq\_P\_TM}
and \texttt{PartIIUniform.NP\_disc\_eq\_NP\_TM}.
Neither theorem establishes
a polynomial-resource physical
realization by the curvature PDE.
The original two-phase construction's
specific uniform writers remain
Restricted in Part II (II-10).
\emph{Addendum, 12 November 2026:} both
writers are now verified as polynomial-time
single-tape transducers, and II-10 is Exact.

\paragraph{Field transport and well-posedness.}
The transport lower bound is verified
in Lean \emph{from} the declared
instantaneous flux inequality; that
implication does not by itself
derive the flux inequality for every
continuum solution.
For the stated finite torus
method-of-lines system, Lean derives
mass and Hamiltonian conservation
and a finite-dimensional flux
inequality with explicit assumptions
(\texttt{PVNP.WCTLattice}).
For a specified frozen cubic--quintic
PDE, Phase VI supplies scoped
$H^2$ periodic-torus well-posedness
and real-detector flux results
(\texttt{PVNP.Torus}).
Those results do not establish
a universal device-size-independent
physical computation rail.
The independent fixed-mass nodal
supplement has an exact finite
two-node Lean analogue and a
separate \emph{analytic} continuum
compactness proof, not a
kernel-formalized continuum theorem.

\paragraph{Open physical and complexity boundaries.}
Unrestricted encoders or decoders
trivialize physical decision models.
The measured physical preparation,
bandwidth, energy, spatial
precision, and readout costs
must be specified and charged.
Neither the finite curvature
inequalities nor the restricted
quadratic single-tape SAT lower
bounds imply a universal
superpolynomial SAT decision bound.
The all-algorithm separator and
locked-expander bounds remain
unproved.  Classical and physical
P versus NP remain unresolved.

\paragraph{Verification and publication provenance.}
Some Phase XII--XIII modules were
merged after GitHub Actions jobs
failed before their first step;
newer source code is not a
freshly passed kernel and axiom
audit.  The exact-head Lean checks
must succeed before a release
may call those declarations
independently verified.
Older statuses remain in Git
history and the 91-claim ledger.
Recompiled PDFs are needed
after the latest source amendments.

\clearpage
\appendix
\section{Conditional Curvature--Capacity Analysis}
\label{sec:curvature_capacity}
\subsection{Finite first-derivative energy is insufficient}
\begin{proposition}[A high-frequency counterexample]
\label{prop:highfreq}
On the periodic cube $\Omega=(0,2\pi)^3$, let
\[
\psi_N(x)=1+\frac1N\sin(Nx_1),\qquad N\ge2.
\]
Then $\int_\Omega|\nabla\psi_N|^2dx=4\pi^3$ and
$\int_\Omega|\Delta\psi_N|^2dx=4\pi^3N^2$.
With $\varepsilon=1$ in Definition~\ref{def:regularized},
\[
C_{\Theta,1}[\psi_N]\ge\frac{16\pi^3}{25}N^2\longrightarrow\infty,
\]
while the ordinary nonnegative mass, quartic, and gradient contributions
of a fixed-coefficient first-derivative field-energy functional are
uniformly bounded.
\end{proposition}
\begin{proof}
$\nabla\psi_N=(\cos(Nx_1),0,0)$ and
$\Delta\psi_N=-N\sin(Nx_1)$, giving the stated integrals. Since
$1/2\le\psi_N\le3/2$, the elementary inequality
$\psi_N/(1+\psi_N^2)\ge2/5$ yields
$|\Theta_1[\psi_N]|\ge(2/5)|\Delta\psi_N|$.
Integrating gives the claimed bound. The amplitude stays between $1/2$
and $3/2$, so its $L^2$ and $L^4$ contributions remain bounded.
\end{proof}
\begin{remark}[WIN-14]
Finite energy involving $\|\nabla\psi\|_2^2$ does not imply a
uniform $H^2$ estimate, and generic finite action does not silently
supply it. A positive $H^2$ regularity assumption or a proved PDE
estimate is required.
\end{remark}

\subsection{Locality and gravitational considerations}
Lieb--Robinson bounds concern information propagation in suitable
short-range quantum lattice systems; they do not imply that arbitrary
PDE second derivatives grow at most polynomially in time. The original
argument also identified divergence of a chosen normalized Laplacian
functional with divergence of the stress--energy tensor without an
established inequality between them. Einstein's equation by itself
cannot close that missing step. The high-frequency family above is a
counterexample to a universal fixed-energy bound on the chosen curvature
functional.

\begin{theorem}[Corrected conditional capacity law]
\label{thm:capacity}
Assume a uniform family of WCT fields $\psi_{n,x}(t)\in H^2(\Omega_n)$
with a fixed $\varepsilon>0$ and
\[
\sup_{|x|=n}\sup_{t\le T(n)}
\|\Delta\psi_{n,x}(t)\|_{L^2(\Omega_n)}\le q(n)
\]
for one polynomial $q$. Then the regularized curvature cost obeys
\[
\sup_{|x|=n}\sup_{t\le T(n)}
\CTheta[\psi_{n,x}(t)]\le\frac{q(n)^2}{4\varepsilon^2}.
\]
\end{theorem}
\begin{proof}
Apply Lemma~\ref{lem:h2bound} uniformly over the family.
\end{proof}
\begin{remark}[WIN-15]
This theorem is conditional on the missing uniform second-derivative
estimate. Finite energy, spatial locality, and finite runtime alone do not
provide the hypothesis. No universal QFT/GR/Lieb--Robinson polynomial
curvature-capacity law is claimed.
\end{remark}

\section{Boolean Influence and Hypercube Boundary}
\label{sec:boundary_capacity}
For a Boolean function $f:\{0,1\}^n\to\{0,1\}$, let $x^{\oplus i}$
denote bit $i$ flipped. Define
\[
\operatorname{Inf}_i(f)=2^{-n}\#\{x:f(x)\ne f(x^{\oplus i})\},
\qquad \operatorname{Var}(f)=\sum_{i=1}^n\operatorname{Inf}_i(f).
\]
Define the \emph{directed} boundary
$B_{\to}(f)=\#\{(x,i):f(x)\ne f(x^{\oplus i})\}$ and the
\emph{undirected} boundary
$B_{\rm und}(f)=\tfrac12B_{\to}(f)$.

\begin{lemma}[Corrected boundary--variation identities]
\label{lem:boundary_variation}
For every Boolean $f$,
\[
B_{\to}(f)=2^n\operatorname{Var}(f),\qquad
B_{\rm und}(f)=2^{n-1}\operatorname{Var}(f).
\]
\end{lemma}
\begin{proof}
Count the $2^n\operatorname{Inf}_i(f)$ directed edges in coordinate
$i$ and sum over $i$; every undirected edge has two orientations.
\end{proof}
\begin{remark}[WIN-20]
The original definition counted directed pairs while its lemma used the
undirected normalization. Both formulas are now displayed separately.
\end{remark}

\begin{theorem}[Universal total-influence bound]
\label{thm:influence}
For any Boolean function of $n$ bits,
\[
0\le\operatorname{Var}(f)\le n.
\]
Consequently, it is impossible to have
$\operatorname{Var}(F_n)=2^{\Omega(n)}$ for any family of Boolean
functions $F_n$ with $n$ variables.
\end{theorem}
\begin{proof}
Each $\operatorname{Inf}_i(f)$ is a probability in $[0,1]$.
Summing $n$ terms proves the upper bound.
\end{proof}

\begin{corollary}[Polynomial variation is vacuous as a complexity class]
\label{cor:var}
If $P_{\rm var}$ denotes Boolean language families with
$\operatorname{Var}(f_n)\le n^{O(1)}$ for all $n$, then
$P_{\rm var}$ contains every Boolean language family. Thus the proposed
exponential-variation exclusion of SAT and the resulting argument for
$P\ne NP$ do not hold.
\end{corollary}
\begin{proof}
Apply Theorem~\ref{thm:influence} with exponent one.
\end{proof}
\begin{remark}[WIN-21]
A large \emph{unnormalized} cube boundary can be exponential even for
simple functions; for example parity has
$\operatorname{Var}(f)=n$ and $B_{\to}(f)=n2^n$.
Neither an exponential directed boundary nor graph expansion by itself
implies classical algorithmic hardness. The previous theorems claiming an
exponential total influence and a strict subclass of $NP$ have been
replaced by these exact bounds.
\end{remark}

\section{Arrow of Time: Limits of the Claim}
The thermodynamic arrow and the availability of time-reversing or
closed-timelike-curve models are additional physical hypotheses, not
logical characterizations of classical P and NP. A process computing a
witness need not run the verifier backward: search algorithms operate
forward in ordinary time. Entropy-reducing maps need not send information
from the future to the past, and inverse dynamics need not efficiently
solve SAT. Conversely, a generic assertion of retrocausality does not
specify the special self-consistent fixed-point model considered in
\cite{AaronsonWatrous2009,Deutsch1991}.

\begin{remark}
No equivalence
$P^{\mathrm{phys}}=NP^{\mathrm{phys}}\iff\text{retrocausality}$,
and no physical separation inferred merely from a forward thermodynamic
arrow, is proved. This section is retained as a research direction rather
than as a theorem.
\end{remark}
\begin{remark}[Addendum, 12 November 2026]
For the frozen dynamics the sentence on inverse dynamics has an exact
form.  Complex conjugation reverses time: if $\psi(t)$ is a solution,
so is $\overline{\psi(-t)}$.  Hence the state at time $-t$ is the
conjugate of the forward evolution, for time $t$, of the conjugated
initial state
(\nolinkurl{Torus.backward_eq_forward_conj}, \nolinkurl{Torus.flow_neg}).
Running this dynamics backward is a forward computation with the
encoder and decoder composed with conjugation.  The remark above remains
an interpretation, not a theorem.
\end{remark}

\section{Mathematical Revision Register}
\label{app:register}
The following identifiers refer to the issue windows attached to this
revised manuscript. ``Corrected'' means a valid replacement statement
has been supplied; ``conditional'' means the original conclusion is not
established and a precise additional hypothesis is specified.
\begin{description}[leftmargin=1.5cm,style=nextline]
\item[WIN-12] Zero-curvature counterexample to the distinct-state lemma
(Proposition~\ref{prop:zero_cost}); corrected.
\item[WIN-13] Unsafe denominator $|\psi+\varepsilon|\ge\varepsilon$;
regularization/gap distinction (Definition~\ref{def:regularized}); corrected.
\item[WIN-14] First-derivative energy does not control $H^2$;
explicit counterexample (Proposition~\ref{prop:highfreq}); corrected.
\item[WIN-15] Universal polynomial curvature growth from locality;
replaced by conditional $H^2$ capacity (Theorem~\ref{thm:capacity}).
\item[WIN-16] Impossible assertion physical NP is empty;
verifier inclusion (Theorem~\ref{thm:p-in-np}); corrected.
\item[WIN-17] Classical mathematical ill-posedness;
replaced by a distinction between formal and physical questions
(Section~\ref{sec:two_semantics}); corrected.
\item[WIN-18] Unsupported expander-to-exponential-interface implication;
conditional interface bound (Proposition~\ref{prop:interface}).
\item[WIN-19] Inadmissible high-curvature NP witnesses;
common physical verifier rail specified (Section~\ref{sec:constructibility}).
\item[WIN-20] Directed/undirected cube boundary factor of two;
exact identities (Lemma~\ref{lem:boundary_variation}); corrected.
\item[WIN-21] Exponential Boolean total influence;
universal bound and vacuity of $P_{\rm var}$ (Theorem~\ref{thm:influence});
corrected.
\end{description}

\clearpage

\begin{thebibliography}{99}

\bibitem{Turing1936}
A.~M.~Turing.
\newblock On Computable Numbers, with an Application to the Entscheidungsproblem.
\newblock \emph{Proceedings of the London Mathematical Society}, 42(2):230–265, 1936.

\bibitem{Church1936}
A.~Church.
\newblock An Unsolvable Problem of Elementary Number Theory.
\newblock \emph{American Journal of Mathematics}, 58(2):345–363, 1936.

\bibitem{Landauer1961}
R.~Landauer.
\newblock Irreversibility and Heat Generation in the Computing Process.
\newblock \emph{IBM Journal of Research and Development}, 5(3):183–191, 1961.

\bibitem{Bennett1982}
C.~H.~Bennett.
\newblock The Thermodynamics of Computation—A Review.
\newblock \emph{International Journal of Theoretical Physics}, 21(12):905–940, 1982.

\bibitem{Lloyd2000}
S.~Lloyd.
\newblock Ultimate Physical Limits to Computation.
\newblock \emph{Nature}, 406:1047–1054, 2000.

\bibitem{Bekenstein1981}
J.~D.~Bekenstein.
\newblock Universal Upper Bound on the Entropy-to-Energy Ratio for Bounded Systems.
\newblock \emph{Physical Review D}, 23(2):287–298, 1981.

\bibitem{AaronsonWatrous2009}
S.~Aaronson and J.~Watrous.
\newblock Closed Timelike Curves Make Quantum and Classical Computing Equivalent.
\newblock \emph{Proceedings of the Royal Society A}, 465:631–647, 2009.

\bibitem{Deutsch1991}
D.~Deutsch.
\newblock Quantum Mechanics Near Closed Timelike Lines.
\newblock \emph{Physical Review D}, 44(10):3197–3217, 1991.

\bibitem{Sipser}
M.~Sipser.
\newblock \emph{Introduction to the Theory of Computation}.
\newblock Cengage Learning, 3rd edition, 2012.

\bibitem{Papadimitriou1994}
C.~H.~Papadimitriou.
\newblock \emph{Computational Complexity}.
\newblock Addison-Wesley, 1994.

\bibitem{AroraBarak2009}
S.~Arora and B.~Barak.
\newblock \emph{Computational Complexity: A Modern Approach}.
\newblock Cambridge University Press, 2009.

\bibitem{HarlowOoguri2018}
D.~Harlow and H.~Ooguri.
\newblock Symmetry, Topology, and Quantum Gravity.
\newblock \emph{Communications in Mathematical Physics}, 383, 2021.

\bibitem{LiebRobinson1972}
E.~H.~Lieb and D.~W.~Robinson.
\newblock The Finite Group Velocity of Quantum Spin Systems.
\newblock \emph{Communications in Mathematical Physics}, 28:251–257, 1972.

% ---------------------------------------------------------------
% Added Wave Confinement Theory / WCC P vs NP References
% ---------------------------------------------------------------

\bibitem{Reyes2025_PWCC}
R.~J.~Reyes.
\newblock P vs NP in Curvature-Bounded Wave Computation: A Model-Relative $P_{\mathrm{WCC}} \neq NP_{\mathrm{WCC}}$ Separation.
\newblock Zenodo, 2025.
\newblock \url{https://doi.org/10.5281/zenodo.17743607}.

\bibitem{Reyes2025_DiscreteWCC}
R.~J.~Reyes.
\newblock Discrete Wave-Constrained Computation and Classical Complexity: Turing Equivalence for $\mathbf{P}$ and $\mathbf{NP}$.
\newblock Zenodo, 2025.
\newblock \url{https://doi.org/10.5281/zenodo.17732642}.

\end{thebibliography}
\end{document}